\SourceLeaf{23}
weil ihre Summe $=0$ ist) sein, die nicht $\equiv 0$ ist. Ebenso wenig
können aber die Größen $\alpha, \beta, \gamma$ alle zueinander inkongruent sein.
Sonst müssten sie, gleichgültig in welcher Reihenfolge, zu den Zahlen
$0, 1, 2$ kongruent und ihre Summe daher durch $3$ teilbar sein. Folglich
ist unter den Differenzen $\alpha-\beta$,
$\beta-\gamma$, $\gamma-\alpha$ genau eine, die
$\equiv 0$ ist. Es gibt somit drei Fälle, die sich gegenseitig ausschließen:

\[
\text{I.} \quad \text{Wenn } \beta \equiv \gamma \quad(\mathrm{mod.}\ 3)
\]
\[
\text{II.} \quad \text{Wenn } \gamma \equiv \alpha \quad(\mathrm{mod.}\ 3)
\]
\[
\text{III.} \quad \text{Wenn } \alpha \equiv \beta \quad(\mathrm{mod.}\ 3).
\]

Welcher dieser drei Fälle eintritt, hängt von $x$ ab. Wir teilen daher
alle durch $p$ nicht teilbaren Zahlen nach diesen drei Fällen in drei
Klassen ein. In die erste kommen diejenigen, bei denen der erste Fall eintritt,
in die zweite diejenigen, bei denen der zweite eintritt, und in die dritte
diejenigen, bei denen der dritte eintritt.

\textls[100]{Erster Satz:} Alle modulo $p$ kongruenten Zahlen gehören
derselben Klasse an.

\textls[100]{Zweiter Satz:} Zahlen mit gleichem Betrag und entgegengesetzten
Vorzeichen gehören derselben Klasse an.

Daraus ist klar, dass man nur die Zahlen $1, 2, 3, \ldots,
\frac{1}{2}(p-1)$
klassifizieren muss, um sofort alle übrigen klassifizieren zu können.

\textls[100]{Dritter Satz:} Die Zahlen $(3n-m)x$ und $2mx$ gehören zwei
aufeinanderfolgenden Klassen an.

\textls[100]{Vierter Satz:} Gilt also $\dfrac{2n-m}{3m} \equiv f$ und damit

\[
1+f+ff \equiv 0 \quad(\mathrm{mod.}\ p), \qquad 1 \equiv f^3 \quad(\mathrm{mod.}\ p),
\]
\noindent so gehören $x$ und die Reste von $fx$ und $ffx$
drei in umgekehrter Reihenfolge aufeinanderfolgenden Klassen an.

